% The bring-up order, and what each wrong order produces.
\begin{tikzpicture}[font=\sffamily\scriptsize, x=1cm, y=1cm]

\tikzset{
  bstep/.style={draw=linecol, fill=kercol, rounded corners=2pt, align=center,
                text width=30mm, minimum height=9mm, inner sep=3pt,
                font=\sffamily\scriptsize},
  bfail/.style={draw=red!70!black, fill=red!8, rounded corners=2pt, align=left,
                text width=46mm, inner sep=3pt, font=\sffamily\scriptsize},
}

\foreach \i/\t in {0/{1. raise the core voltage\\{\tiny voltage scale first}},
                   1/{2. set the flash wait states\\{\tiny before any speed increase}},
                   2/{3. start the phase-locked loop\\{\tiny 8 MHz / 2, times 140, / 2}},
                   3/{4. switch the system clock\\{\tiny and only now}},
                   4/{5. select the peripheral kernel clocks\\{\tiny one register per group}},
                   5/{6. read every frequency back\\{\tiny from the registers, not the arithmetic}}}{
  \node[bstep] (s\i) at (0,-\i*1.95) {\t};
}
\foreach \i in {0,1,2,3,4}{ \draw[flow] (s\i) -- (s\the\numexpr\i+1\relax); }

% ---------------- what each wrong order produces
\node[bfail] (f0) at (5.4,0)      {Skip it and the part runs below its rated frequency, or stops at the phase-locked loop: the voltage scale is what permits 280 MHz at all.};
\node[bfail] (f1) at (5.4,-1.95)  {Skip it and the core fetches from flash faster than the flash can answer. The failure is immediate and looks like a hardware fault.};
\node[bfail] (f2) at (5.4,-3.90)  {Take the multiplier and dividers from material written for the sibling part, and the tree does not lock or locks at the wrong frequency.};
\node[bfail] (f3) at (5.4,-5.85)  {Switch before the loop is locked and the part halts. Switch and forget to check, and the node runs at reset speed with every later timing figure wrong.};
\node[bfail] (f4) at (5.4,-7.80)  {Use the sibling's register name and the build fails, which is the good case. Use the right register with the wrong field and a peripheral is clocked from something unintended.};
\node[bfail] (f5) at (5.4,-9.75)  {Skip it and every derived value in chapters 2 and 3 rests on arithmetic nobody checked.};

\foreach \i in {0,1,2,3,4,5}{ \draw[faultedge] (s\i.east) -- (f\i.west); }

\node[chlabel, anchor=north west, text width=44mm] at (-1.9,-11.3)
  {This sequence is inherited without thought by every chapter after this one,
   which is exactly why it is written down once, here, with the failure mode of
   each step beside it.};
\node[umlnote, text width=52mm, anchor=north west] at (3.4,-11.3)
  {Step 6 is the one most projects omit. The arithmetic in step 3 is a hope; the
   register read in step 6 is a fact, and the two disagreeing is a board that
   boots and lies about its own speed. This node refuses to start when they
   disagree.};
\end{tikzpicture}
